\section*{Erd\H{o}s Problem 1138: a prime-gap obstruction}

\paragraph{Statement and scope.}
Let \(p_n\) denote the increasing primes and
\(d(x)=\max_{p_n<x}(p_{n+1}-p_n)\).
The proposed assertion is that for every fixed \(C>1\), uniformly over \(x/2<y<x\),
\[
 \pi(y+Cd(x))-\pi(y)\sim \frac{Cd(x)}{\log y}.
\]
We disprove this universal assertion. The proof reconstructs the record-gap argument reported at \url{https://www.erdosproblems.com/1138} and in the linked Sunder--Kumrawat--Cheri write-up, which credits GPT-5.5 assistance:
\url{https://sourish-kumrawat.github.io/papers/Erdos_1138.pdf}.

\paragraph{Record gaps.}
Prime gaps are unbounded, since factorial blocks give arbitrarily long strings of composite integers. Consequently there are infinitely many increasing record gaps \(d=q-p\), where \(p,q\) are consecutive primes and no earlier prime gap is larger. The prime number theorem implies \(d=o(p)\) along these records: for each fixed \(\varepsilon>0\) there is a prime in \((p,(1+\varepsilon)p)\) for all sufficiently large \(p\).

Set \(x=p+1\) and \(y=p-2d\). Then \(d(x)=d\); this includes the gap beginning at \(p\), and there is no further prime below \(x\). For sufficiently large records, \(x/2<y<x\).
The two endpoints
\[
 y+2d=p,\qquad y+\tfrac52d=p+\tfrac12d
\]
have the same prime counting function, because \((p,q)\) contains no prime. Hence the two corresponding numerators have a common value \(M\).

\paragraph{Incompatible asymptotics.}
Put \(T=d/\log y>0\). The proposed asymptotics for the two fixed constants \(C=2\) and \(C=5/2\) would require simultaneously
\[
 M/(2T)\longrightarrow1,\qquad M/((5/2)T)\longrightarrow1,
\]
which is impossible. More quantitatively, for every \(M\geq0\),
\[
 \max\left\{\left|\frac{M}{2T}-1\right|,
              \left|\frac{M}{(5/2)T}-1\right|\right\}\geq\frac19.
\]
Otherwise the two allowed intervals for \(M/T\), whose common boundary would be \(20/9\), are disjoint. One of these two fixed constants therefore fails along infinitely many of the records, by the infinite pigeonhole principle.

This refutes the assertion quantified over every \(C>1\). It does not claim that the displayed argument singles out which of \(2\) and \(5/2\) fails, or that every fixed \(C\) fails.

\paragraph{Completion estimate.}
The full universal assertion is disproved, using the prime number theorem as an explicit classical input.
\textbf{COMPLETION: 100\%}
